\chapter{Introducción y puesta en marcha}
\label{ch:intro}

\section{Qué es Fitbauer}

\textbf{Fitbauer} es una aplicación de escritorio para el análisis de espectros
Mössbauer de $^{57}$Fe. Permite:

\begin{itemize}[nosep]
  \item \textbf{Cargar} espectros en varios formatos (cuentas en bruto
        \ruta{.ws5}/\ruta{.adt} o espectros ya doblados en velocidad
        \ruta{.csv}/\ruta{.txt}/\ruta{.dat}/\ruta{.exp}).
  \item \textbf{Doblar} (\emph{folding}) las cuentas en bruto y construir el eje
        de velocidad a partir de una \emph{calibración}.
  \item \textbf{Simular} el modelo directo con hasta tres componentes discretos
        (singlete, doblete, sexteto) y perfiles Lorentz/Voigt.
  \item \textbf{Ajustar} de forma discreta o mediante \emph{distribuciones}
        $P(\BHF)$ / $P(\dEQ)$.
  \item \textbf{Documentar}: informes Markdown/PDF, exportación de gráficos
        interactivos y sesiones reproducibles.
\end{itemize}

La interfaz está disponible en español, inglés y francés (menú \menu{Idioma}).

\begin{nota}
El cálculo (física y motores de ajuste) reside en el paquete \ruta{core/}, que
no depende de ninguna interfaz gráfica. La GUI es un cliente de ese núcleo. Este
manual describe el uso de la aplicación; el detalle matemático del motor de
ajuste está en el documento aparte \emph{Manual matemático del motor de ajuste}.
\end{nota}

\section{Instalación y arranque}

Fitbauer requiere Python 3.11 o superior. Desde la raíz del repositorio:

\begin{lstlisting}[language=bash]
python3 -m venv .venv
source .venv/bin/activate        # Windows: .venv\Scripts\activate
pip install -r requirements.txt
python fitbauer.py               # abre la interfaz gráfica
\end{lstlisting}

El punto de entrada único es \ruta{fitbauer.py}, que lanza la ventana principal
(Qt + Matplotlib). También existen versiones ejecutables empaquetadas
(PyInstaller) para quienes no deseen instalar Python.

\section{Panorama de la interfaz}
\label{sec:panorama}

Al arrancar aparece la ventana principal, organizada en:

\begin{itemize}[nosep]
  \item \textbf{Barra de menús}: \menu{Archivo}, \menu{Ajuste}, \menu{Vista},
        \menu{Idioma} y \menu{Ayuda}.
  \item \textbf{Panel de calibración}: control de \param{Vmax}, centro de
        \emph{folding}, línea base, pendiente y perfil de línea.
  \item \textbf{Paneles de componentes} (hasta tres): forma del componente y sus
        parámetros físicos ($\delta$, $\dEQ$, $\BHF$, anchuras, intensidades).
  \item \textbf{Panel de distribución}: modo $P(\BHF)$/$P(\dEQ)$, forma,
        regularizador y sus parámetros.
  \item \textbf{Zona de gráfico}: representación del espectro, el modelo y los
        subespectros, con barra de herramientas de navegación (Matplotlib).
  \item \textbf{Panel de información/estado}: parámetros ajustados, estadísticos
        ($\chi^2$, AIC, BIC) y mensajes.
\end{itemize}

\guicap{gui_overview}{Vista general de la ventana principal de Fitbauer con sus
paneles. (Sustituir por una captura real; nombre del fichero indicado bajo el
recuadro.)}

\begin{truco}
La disposición de los paneles es configurable desde
\menu{Vista > Configurar layout de paneles…}. Puede guardar presets adaptados a
pantallas grandes o portátiles.
\end{truco}

\section{El flujo de trabajo típico}

El orden recomendado, que estructura los capítulos siguientes, es:

\begin{enumerate}[nosep]
  \item \textbf{Calibrar} (capítulo~\ref{ch:calibracion}): cargar o definir una
        calibración de velocidad y \textbf{fijarla}. A partir de ese momento se
        reutiliza automáticamente.
  \item \textbf{Cargar datos} (capítulo~\ref{ch:carga}): abrir el espectro
        problema; la calibración fijada se aplica sola.
  \item \textbf{Simular y ajustar} (capítulo~\ref{ch:simular}): construir el
        modelo directo y ajustarlo; validar con un calibrado de $\alpha$-Fe.
  \item \textbf{Distribuciones} (capítulo~\ref{ch:distribuciones}): cuando el
        campo hiperfino no es único, simular y ajustar $P(\BHF)$.
\end{enumerate}
